\section{测度的收敛与极限操作}
\label{sec:measure-convergence}

本节把前面的比较量用于测度序列和时间参数曲线。
先区分收敛方式，再建立弱极限下事件、映射和成本的可用判据，随后说明什么时候存在极限子列，最后简述时间参数曲线。
各项分别承担定义、稳定操作、极限存在和应用的职责，不把它们混作一种收敛分类。
比较始终发生在固定空间中；变空间对象先按上一节给出共同表示。

\subsection{测度序列的收敛方式}
\label{sec:measure-sequence-convergence}

同一概率列可以由不同观测判定是否趋近：全部事件的统一误差、有界连续观测量，或者加入距离矩的尾部控制。
每种方式的对象、量词和所需附加结构都要分别写清。

\subsubsection{\term{逐集收敛}与全变差收敛}

\paragraph{从比较误差定义趋近}

固定概率列$P_n$和候选概率$P$。
最强的常用事件判据是
\[
 \operatorname{TV}(P_n,P)\to0:
\]
所有可测事件的误差用同一个趋于零的量统一控制。
由\ref{sec:measure-metric-relations}小节的双向不等式，这等价于$H(P_n,P)\to0$；$D_{\mathrm{KL}}(P_n\|P)\to0$是取得这种控制的充分保证（Pinsker界），但由$P_\varepsilon=(1-\varepsilon)\delta_0+\varepsilon\delta_1$与$\delta_0$的例子可知不能反推。

更弱的事件判据是逐集收敛：对每个固定事件$A\in\mathcal A$，
\[
 P_n(A)\to P(A).
\]
给定误差精度后，起始指标$N$允许依赖事件$A$；全变差收敛则要求同一个$N$对全部事件有效。
所以全变差收敛推出逐集收敛，反之一般不成立。
对固定的有界可测函数$g$，逐集收敛仍能控制其积分：先用有限简单函数在$\|g\|_\infty$意义下一致逼近$g$，再用概率总质量为$1$控制两端逼近误差，得到
\[
 \int g\,dP_n\to\int g\,dP.
\]

\subsubsection{\term{弱收敛}与连续测试}

\paragraph{弱收敛与位置移动}

在固定度量空间的Borel概率上，定义$P_n$弱收敛到$P$，记为$P_n\Rightarrow P$，若对每个有界连续函数$\varphi$都有
\[
 \int\varphi\,dP_n\to\int\varphi\,dP.
\]
测试类由“全部可测事件”缩小为“有界连续观测量”。
由于事件指标一般不连续，逐集收敛能推出弱收敛（对每个固定有界连续函数使用上面的积分结论），但弱收敛不自动控制不连续指标，也不自动控制无界矩。

在实轴上取
\[
 P_n=\delta_{1/n},\qquad P=\delta_0.
\]
对任意有界连续$\varphi$，测试积分就是$\varphi(1/n)\to\varphi(0)$，故$P_n\Rightarrow P$；但固定事件$\{0\}$的质量差始终为$1$，\term{全变差距离}也始终为$1$。
同一个位置移动的例子说明，采用哪种比较标准会改变“趋近”的结论。

\subsubsection{Wasserstein收敛与矩条件}

\paragraph{弱趋近与距离尾部的共同控制}

在本章的\term{Polish空间}范围内，$1\leq p<\infty$，$P_n,P\in\mathcal P_p(X)$，Wasserstein收敛有标准判据：
\begin{equation}
 W_p(P_n,P)\to0
 \quad\Longleftrightarrow\quad
 P_n\Rightarrow P
 \ \text{且}\
 \int d(x,x_0)^p\,dP_n\to\int d(x,x_0)^p\,dP.
 \label{eq:measure-wasserstein-criterion}
\end{equation}
这里把判据作为标准结论使用，下面解释矩条件限制的究竟是什么。

令$z(x)=d(x,x_0)^p$。
弱收敛和全矩收敛成立时，对有界连续截断$\min(z,R)$先用弱收敛：
\[
 \int\min(z,R)\,dP_n\to\int\min(z,R)\,dP.
\]
再从全矩中减去有限截断部分，得到
\[
 \int(z-R)_+\,dP_n\to\int(z-R)_+\,dP.
\]
先取足够大的$R$使极限尾部小，再让$n$足够大，并用有限个首项各自可积处理剩余项，就能取得对全列的共同尾部控制。
利用
\[
 z\,\mathbf 1_{\{z>2R\}}\leq2(z-R)_+,
\]
共同超额控制转化为真正的尾积分控制。
反向则由截断积分收敛和共同尾部控制恢复全矩收敛。
所以弱收敛之外还需要距离尾部的一致可积性，而不只是每一项矩分别有限。

\paragraph{小概率远距离的反例}

在实轴上取$n\geq2$及
\[
 P_n=(1-n^{-p})\delta_0+n^{-p}\delta_n,\qquad P=\delta_0.
\]
对任意有界连续$\varphi$，远端项的贡献被$n^{-p}\|\varphi\|_\infty$控制并趋于零，故$P_n\Rightarrow P$；每一项也都有有限$p$阶矩。
但把远端质量搬回零点的成本为
\[
 W_p(P_n,P)^p=n^{-p}\cdot n^p=1,
\]
不趋于零。
这个反例把“远端质量很少”与“搬运成本很小”分开：缺少共同的距离尾部控制时，小概率乘大距离仍可保持常数。
涉及映射后的分布时同样要先检查输出矩，不能只因输入矩有限就使用目标空间的Wasserstein判据。

\subsection{弱收敛的判据与稳定性}

弱收敛定义只直接控制有界连续测试。
本节分别问事件概率怎样读取、连续变换后的测度是否仍收敛、无界成本能保留什么不等式。
这些是同一极限方式下的可用操作，不是与弱收敛并列的新收敛类型。

\subsubsection{事件概率与边界质量}

在Polish空间的Borel概率范围内，\term{Portmanteau定理}给出常用等价形式：
\[
 P_n\Rightarrow P
\]
等价于对所有闭集$C$有
\[
 \limsup_{n}P_n(C)\leq P(C),
\]
也等价于对所有开集$G$有
\[
 \liminf_nP_n(G)\geq P(G).
\]
由此，若Borel集$A$满足边界零质量$P(\partial A)=0$，则
\[
 P_n(A)\to P(A).
\]
这里只展开“从弱收敛推出这些可用界”的方向，整套等价形式作为标准定理使用。

对非空闭集$C$，取
\[
 h_k(x)=\max\{1-k\,d(x,C),0\}.
\]
每个$h_k$有界连续，且随$k$递减到$\mathbf 1_C$。
固定$k$时，由$\mathbf 1_C\leq h_k$和弱收敛，
\[
 \limsup_nP_n(C)\leq\lim_n\int h_k\,dP_n=\int h_k\,dP.
\]
再对固定极限测度用支配收敛令$k\to\infty$，得到闭集上界$\limsup_nP_n(C)\leq P(C)$；空闭集单独为零。
开集下界由补集得到：
\[
 \liminf_nP_n(G)\geq P(G).
\]
对一般集合$A$，用内部$A^\circ$和闭包$\overline A$夹住：
\[
 P(A^\circ)\leq\liminf_nP_n(A)\leq\limsup_nP_n(A)\leq P(\overline A).
\]
只有边界质量为零时两端才相等。
回看$\delta_{1/n}$与事件$\{0\}$：极限测度在边界$\{0\}$上有质量$1$，不满足条件，所以$P_n(\{0\})$不收敛到$P(\{0\})$。
事件收敛的例外出现在测试的不连续位置，不能只用“弱收敛比较宽松”一句话带过。

\subsubsection{连续推前与边缘保持}

设$T:X\to Y$是固定连续映射，$P_n\Rightarrow P$。
对任意目标端有界连续函数$\varphi$，先用推前积分公式改写：
\[
 \int_Y\varphi\,d(T_\#P_n)=\int_X\varphi\circ T\,dP_n.
\]
复合$\varphi\circ T$仍有界连续，由弱收敛右端趋于$\int_X\varphi\circ T\,dP$，所以
\[
 T_\#P_n\Rightarrow T_\#P.
\]
连续映射保持弱极限；这一步明确了在哪个空间用哪类测试，不能把结论推广到任意\term{可测映射}。

可测但不连续的映射可能失败。
在实轴上取$P_n=\delta_{1/n}$、$P=\delta_0$和$T(x)=\mathbf 1_{\{0\}}(x)$，则$T_\#P_n=\delta_0$而$T_\#P=\delta_1$，并不弱收敛。
连续坐标投影是乘积空间上的连续映射，所以弱收敛\term{耦合}列的边缘继承相应边缘极限；这是后面取得最优计划时保持边缘约束的依据。

\subsubsection{\term{下半连续}成本的极限下界}

对非负下半连续函数$c:X\to[0,\infty]$，可以取非负有界连续函数从下方逼近：
\[
 c_k(x)=\min\left\{k,\ \inf_y[c(y)+k\,d(x,y)]\right\}.
\]
每个$c_k$有界且连续（Lipschitz常数为$k$），随$k$递增并不超过$c$。
若在某点$x$有$\sup_kc_k(x)<c(x)$，取近似极小点$y_k$，则$d(x,y_k)\to0$而$y_k\to x$，下半连续性给出$c(x)\leq\liminf_kc(y_k)\leq\sup_kc_k(x)$，矛盾；故$c_k\uparrow c$。
若$c$处处为无穷，直接取$c_k=k$即可。

现在设$P_n\Rightarrow P$。先固定$k$，由$c_k\leq c$和弱收敛，
\[
 \int c_k\,dP=\lim_n\int c_k\,dP_n\leq\liminf_n\int c\,dP_n.
\]
再对极限测度使用单调收敛并令$k\to\infty$，得到下半连续积分界
\begin{equation}
 \int c\,dP\leq\liminf_n\int c\,dP_n.
 \label{eq:measure-lsc-portmanteau}
\end{equation}
这里测度随$n$变化，不能直接用固定测度的Fatou不等式替代整个论证；允许成本无界也不意味着积分取等号。

这条下界正是运输存在性论证中的成本步骤：接近最优的耦合列取弱子列后，极限成本不超过原列成本的下极限。
紧子列条件在下一小节核对；非空、有限和最小值达到仍是三个不同判断。

\subsection{\term{一致紧性}与弱子列}
\label{sec:measure-tightness}

前面先给定候选极限，再判断某种收敛是否成立。
本小节反过来问：一族概率是否至少留下一个可用的极限子列？
关键是能否用共同紧集把所有测度的大部分质量留在有界范围内；最后把这一判断用于\term{运输计划}的极小化序列。

\paragraph{一致紧性与极限子列}

对概率族$\mathcal F$定义一致紧性：对每个$\varepsilon>0$，都存在同一个紧集$K_\varepsilon$，使所有$P\in\mathcal F$满足
\[
 P(X\setminus K_\varepsilon)<\varepsilon.
\]
紧集可以随精度变化，但不能随族中的测度变化。
在无限维或一般度量空间中，有界球不必紧，统一矩界也不能无条件替代这一要求。

\begin{theorem}[\term{Prokhorov定理}，子列形式]
\label{thm:measure-prokhorov}
在Polish空间上，概率族$\mathcal F$一致紧，当且仅当$\mathcal F$中的每个序列都有弱收敛到某个Borel概率的子列。
\end{theorem}

极限概率不必属于原来的族$\mathcal F$；“定义某种收敛方式”与“保证存在极限子列”是两个判断。
这里说明使用条件，不展开完整证明。
实轴上的族$\{\delta_n:n\geq1\}$给出质量逃逸的典型反例：每个单项都有紧支撑$\{n\}$，但任何共同紧集都无法留住全部质量，序列没有弱收敛到概率的子列。
它与$P_n=(1-n^{-p})\delta_0+n^{-p}\delta_n$的矩尾反例分别说明“存在弱子列”与“几何成本趋零”需要不同控制，不能混为同一条件。

\paragraph{固定边缘下的最优计划}

在两个Polish空间上固定Borel概率$P,Q$。
先说明整个耦合族$\Pi(P,Q)$一致紧。
给定$\varepsilon>0$，由$P,Q$各自一致紧（单个概率在Polish空间上必一致紧），取紧集$K,L$使
\[
 P(K^c)<\frac\varepsilon2,\qquad Q(L^c)<\frac\varepsilon2.
\]
对任意$\gamma\in\Pi(P,Q)$，边缘约束给出
\[
 \gamma((K\times L)^c)
  \leq\gamma(K^c\times Y)+\gamma(X\times L^c)
  =P(K^c)+Q(L^c)<\varepsilon.
\]
所以是整族耦合共用同一个紧集$K\times L$，不是只证明某个候选单独紧。

设成本$c:X\times Y\to[0,\infty]$非负下半连续，并且至少有一个有限成本耦合。
令最优值
\[
 I=\inf_{\gamma\in\Pi(P,Q)}\int c\,d\gamma<\infty.
\]
按下确界选取耦合列$\gamma_n$使$\int c\,d\gamma_n\to I$。
由耦合族一致紧和Prokhorov定理，取弱收敛子列$\gamma_{n_k}\Rightarrow\gamma$。
连续坐标投影保持弱极限，所以$\gamma$的两个边缘仍为$P,Q$，即$\gamma\in\Pi(P,Q)$。
再由下半连续积分界~\eqref{eq:measure-lsc-portmanteau}，
\[
 \int c\,d\gamma\leq\liminf_k\int c\,d\gamma_{n_k}=I.
\]
而$\gamma$可行又给出$\int c\,d\gamma\geq I$，故最小值在$\gamma$处达到。
这就把紧性、边缘保持和成本下界三步接齐。
结论不保证最优确定性映射存在；下半连续成本的论证也不能用于任意可测成本。

\subsection{测度曲线与参数极限}

\paragraph{由序列判据读取时间连续性}

对固定空间上的概率曲线$t\mapsto\mu_t$，称它弱连续，若每个有界连续测试的读数$t\mapsto\int\varphi\,d\mu_t$连续；在时间区间上等价于
\[
 t_n\to t\ \Longrightarrow\ \mu_{t_n}\Rightarrow\mu_t.
\]
讨论$W_p$连续性时限定各时刻$\mu_t\in\mathcal P_p(X)$，并要求
\[
 W_p(\mu_{t_n},\mu_t)\to0,
\]
由矩判据~\eqref{eq:measure-wasserstein-criterion}，这还同时控制距离尾部。

在实轴上取$\mu_t=\delta_t$，则
\[
 W_p(\mu_s,\mu_t)=|s-t|,
\]
曲线关于$W_p$连续；但$s\neq t$时$\operatorname{TV}(\mu_s,\mu_t)=1$，曲线在全变差意义下处处不连续。
曲线是否连续取决于检测方式。
若每个时刻的测度原本位于不同空间，必须先固定到共同空间的表示，不能在不同底集上直接沿用同一个测试函数或距离。

\paragraph{固定时刻、有限时域与长时间}

对目标曲线$\mu_t$和近似曲线$\mu_t^h$，要区分三种越来越强的趋同说法：
\begin{enumerate}
 \item 每个固定时刻$t$，误差$W_p(\mu_t^h,\mu_t)\to0$；
 \item 每个有限时间区间$[0,T]$上，$\sup_{0\leq t\leq T}W_p(\mu_t^h,\mu_t)\to0$；
 \item $t\to\infty$时，曲线本身是否存在时间极限。
\end{enumerate}
三种量词不能互相替代。

在实轴上取$h>0$、$t\geq0$，令
\[
 \mu_t=\delta_0,\qquad \mu_t^h=\delta_{ht}.
\]
固定$t$时误差为$ht\to0$；有限时域$[0,T]$上的上界为$hT\to0$；但在$t=1/h$这一随$h$移动的时刻误差恒为$1$，不存在关于$h$的全域一致收敛。
这个短例保留了双极限边界。
演化方程如何产生测度曲线、数值近似为何稳定以及何时能交换极限，都留给后续相应问题。
